\documentstyle[%


righttag%


,11pt%


,amssymb%


,russian%


,izvru%


]{amsart}





\newcommand{\tl}[3]{\medskip\noindent{\bf#1}\ #2 \dotfill #3 \par} 


\pagestyle{empty}


\begin{document}


%\tableofcontents


\par


\vskip 2cm


\par


\bigskip


\bigskip


\bigskip


\par


\centerline{\Large СОДЕРЖАНИЕ}


\bigskip


\bigskip


%%%%%%%%%% Use \newline %%%%%%%%%%%





\tl{Аксентьев Л.А.}{Многолистные функции


из расширенных классов Беккера и Нехари\newline


и их гидромеханическое истолкование}{3}


\tl{Александров И.А.}{Конформные отображения полуплоскости 


на области с симметрией\newline переноса}{15}


\tl{Бурмистрова А.Б.}{Относительный коэффициент сюръективности 


линейных операторов}{19}


\tl{Ишанкулов Т.И., Махмудов О.И.}{Задача Коши для системы уравнений 


термоупругости\newline в пространстве}{27}


\tl{Кутрунов В.Н., Курята З.С.}{Интегральные уравнения векторного поля}{33}


\tl{Кытманов А.М., Ходос О.В.}


{Об условиях голоморфного продолжения гладких\newline $CR$-функций


в фиксированную область}{37}


\tl{Лосев А.Г., Мазепа Е.А.}{Об асимптотическом поведении решений некоторых 


уравнений\newline эллиптического типа на некомпактных 


римановых многообразиях}{41}


\tl{Непомнящих Ю.В., Поносов А.В.}{Локальные операторы в некоторых 


подпространст-\newline вах пространства $L_0$}{50}


\tl{Самсонов Б.Ф.}{О преобразовании Дарбу $N$-го порядка}{65}


\tl{Сижук В.П.}{Окрестности функций, выпуклых относительно 


симметричных точек}{69}


\tl{Скимель В.Н.}{К задачам приемлемости приближенных 


решений дифференциальных\newline уравнений}{74}


\tl{Сосов Е.Н.}{Об аппроксимативных свойствах множеств 


в специальных метрических\newline пространствах}{81}











\vfill


\noindent\raisebox{2pt}{\copyright} Казанский госудаpственный унивеpситет


\end{document}


